\section{Introdução}

\begin{frame}{Doença de Parkinson e tremor}
  \begin{columns}[T]
    \begin{column}{0.48\textwidth}
      \begin{block}{A doença}
        \begin{itemize}
          \item Distúrbio neurodegenerativo crônico e progressivo: perda de neurônios dopaminérgicos na substância negra \cite{Meng:2022,Fujikawa:2023}.
          \item Manifestações motoras: bradicinesia, rigidez, instabilidade postural e \alert{tremor de repouso} \cite{Botzel:2014}.
        \end{itemize}
      \end{block}
    \end{column}
    \begin{column}{0.48\textwidth}
      \begin{block}{Tratamentos convencionais}
        \begin{itemize}
          \item Farmacológico: primeira linha, mas parte dos pacientes não responde, e o uso prolongado perde eficácia e causa efeitos colaterais \cite{Fujikawa:2023,Lora-Millan:2021,Meng:2022}.
          \item Cirúrgico: invasivo e com riscos \cite{Botzel:2014,Fujikawa:2023}.
        \end{itemize}
      \end{block}
    \end{column}
  \end{columns}

  \vspace{0.4cm}
  \begin{exampleblock}{Órteses com controle ativo}
    Alternativa eficaz e não invasiva \cite{Lora-Millan:2021,Meng:2022}, particularmente popular para a supressão do tremor de punho.
  \end{exampleblock}
\end{frame}

\begin{frame}{Supressão ativa do tremor de punho}
  \centering
  \small
\begin{tabularx}{\textwidth}{
    l
    >{\raggedright\arraybackslash\hsize=1.2\hsize}X
    >{\raggedright\arraybackslash\hsize=0.8\hsize}X
}
    \toprule
    Trabalho & Controle e planta & Estimação \\
    \midrule
    \textcite{Gallego:2013} & PI (neuroprótese), punho com 1 GdL & CDF + WFLC + KF \\\midrule
    \textcite{Copur:2019} & Controle repetitivo (estimulação elétrica), 1 GdL & Passa-altas de fase zero \\\midrule
    \textcite{Lekshmi:2019} & PID com ganhos por algoritmo genético, braço com 4 GdL (punho com 1 GdL) & Movimento voluntário ausente \\\midrule
    \textcite{Zamanian:2019} & Filtro \textit{notch} adaptativo, punho com 1 GdL & Frequência do tremor via AFE \\\midrule
    \textcite{Qi:2024} & ADRC, punho com 1 GdL & Movimento voluntário via EBMFLC \\
    \bottomrule
\end{tabularx}


  \vspace{0.3cm}
  \begin{alertblock}{Lacuna}
    Tremores de ombro e cotovelo afetam significativamente o movimento do punho,
    mas a maioria das abordagens usa modelos de uma única articulação, desprezando o acoplamento entre elas.
  \end{alertblock}
\end{frame}

\begin{frame}{Proposta deste trabalho}
  \begin{itemize}
    \setlength{\itemsep}{0.25cm}
    \item \textbf{Controle:} rejeição ativa de perturbações baseada em erro (EADRC), robusta a perturbações externas e incertezas de modelo \cite{Guo:2016}.
    \item \textbf{Plataforma de teste:} modelo linear do membro superior com 3 graus de liberdade (GdL) \cite{Gebai:2017,Gebai:2018}, que captura o acoplamento ombro--cotovelo--punho.
    \item \textbf{Estimação do movimento voluntário}, em duas variantes:
      \begin{itemize}
        \item EADRC+EBMFLC, com o combinador linear de Fourier de \textcite{Qi:2024};
        \item EADRC+ZPLP, com filtragem passa-baixas de fase zero (ZPLP).
      \end{itemize}
    \item \textbf{Comparação} com dois PIDs: IMC-PID (\textit{Internal Model Control}) e DE-PID (evolução diferencial com conhecimento perfeito do movimento voluntário).
    \item \textbf{Robustez:} 100 simulações de Monte Carlo dos parâmetros físicos e teste de Wilcoxon pareado.
  \end{itemize}
\end{frame}
